%% kapa: exam=1 round=27 year=2016 period=1 form=A subject=civil_law subject_ko=민법 seq=19 number=19
%% source: Q-Net 감정평가사 기출문제 (cst00309) artlSeq=5169402 file=extracted/1008/5169402/제27회 감정평가사 1차 문제/제27회 감정평가사 1교시문제(A형).pdf
%% conversion: kordoc 4.24.51 -> tools/gen1.py (auto); review=proofread-sonnet; flags=none
\begin{question}{19}
법률행위의 부관에 관한 설명으로 옳지 않은 것은? (다툼이 있으면 판례에 따름)
\begin{choices}
  \item 조건을 붙이고자 하는 의사가 있더라도 그것이 표시되지 않으면 법률행위의 부관으로서의 조건이 되는 것은 아니다.
  \item 어떤 조건이 붙어 있었는지 아닌지는 그 조건의 존재를 주장하는 자가 이를 증명하여야 한다.
  \item 당사자는 조건의 성부가 미정인 동안에 조건의 성취로 인하여 생길 상대방의 이익을 해하지 못한다.
  \item 조건의 내용 자체가 불법적인 것이어서 무효일 경우, 그 법률행위 전부가 무효로 된다.
  \item 부관에 표시된 사실이 발생하지 아니하는 것이 확정된 때에도 그 채무를 이행하여야 한다고 보는 것이 상당한 경우, 조건부 법률행위로 보아야 한다.
\end{choices}
\end{question}
\answer{5}
